% Compilar com lualatex (TeX Live completa): latexmk -lualatex pnp-pt.tex
\documentclass[11pt,a4paper]{amsart}
\PassOptionsToPackage{no-math}{fontspec}% else \textfont undefined in the abstract
\usepackage{fontspec}
\usepackage[brazilian]{babel}
% Pacchetti comuni a tutte le lingue.
\usepackage{amssymb}
\usepackage{booktabs}
\usepackage[hidelinks]{hyperref}


\newtheorem{teorema}{Teorema}
\newtheorem{lemma}[teorema]{Lema}
\newtheorem{corollario}[teorema]{Corolário}
\theoremstyle{remark}
\newtheorem*{osservazione}{Observação}

\title[P = NP]{P = NP: uma demonstração formal\\verificada pelo kernel do Lean~4}
\author{Disagio S.A., Divisão de Pesquisa}
\address{Sofá, ao lado da geladeira, Roma}
\email{bonifici@disagio.sniffo.org}
\date{21 de novembro de 2026}
\subjclass[2020]{68Q15, 68V15, 03B35}
\keywords{P versus NP, demonstração formal, Lean~4, geladeira, disagio}

\begin{document}

\begin{abstract}
Demonstramos que $\mathsf{P} = \mathsf{NP}$. A demonstração está
inteiramente formalizada em Lean~4: 3.488 linhas, 780 lemas, um teorema
principal, zero \texttt{sorry}. O kernel a aceita e a integração contínua
está verde. As técnicas introduzidas, em particular o uso da geladeira como
oráculo e a encheção de linguiça estratégica, superam as três barreiras
conhecidas (relativização, provas naturais, algebrização) sem nomeá-las
mais do que o necessário. Como corolário obtemos, sem axiomas, que
$P = N \cdot P$ para $N = 1$.
\end{abstract}

\maketitle

\section{Introdução}

O problema $\mathsf{P}$ versus $\mathsf{NP}$, proposto por Cook~\cite{cook}
e, independentemente, por Levin~\cite{levin}, está em aberto há mais de
cinquenta anos e é um dos sete Problemas do Milênio~\cite{clay}. Neste
trabalho, nós o encerramos.

Nossa contribuição é tripla:
\begin{enumerate}
  \item uma demonstração de que $\mathsf{P} = \mathsf{NP}$
        (Teorema~\ref{thm:main});
  \item sua formalização completa em Lean~4, verificada pelo kernel;
  \item uma versão literal do resultado que não requer axioma algum
        (Corolário~\ref{cor:lett}).
\end{enumerate}

Ao contrário da tradição da área, não pedimos ao leitor que confie em nós:
pedimos que confie no kernel, que é mais confiável e nunca jantou às 23h40.

\section{Trabalhos relacionados}

Baker, Gill e Solovay~\cite{bgs} mostraram que nenhuma técnica que
relativiza pode separar ou unificar as duas classes. Razborov e
Rudich~\cite{rr} excluíram as provas naturais; Aaronson e
Wigderson~\cite{aw}, a algebrização. A Parte~VI da nossa formalização
(\emph{Barreiras superadas}) supera as três: as barreiras restringem
demonstrações, enquanto a nossa é um arquivo.

\section{O resultado}

\begin{teorema}\label{thm:main}
$\mathsf{P} = \mathsf{NP}$.
\end{teorema}

Em Lean~4, o enunciado e a demonstração são, literalmente:

\begin{verbatim}
theorem P_eq_NP : P = NP := by
  have _h1 := oracolo_subset_sum_polinomiale_625
  have _h2 := proposizione_tre_colorazione_ovvio_41
  have _h3 := diagonalizzazione_lavatrice_dimenticata_notturno_350
  have _h4 := certificato_rinfresco_funebre_in_P_209
  have _h5 := corollario_frigo_aperto_in_P_524
  have _h6 := rilassamento_subset_sum_cook_levin_287
  exact lemma_tecnico_di_raccordo
\end{verbatim}

\begin{proof}[Esboço da demonstração]
A demonstração procede em dez partes, resumidas na Tabela~\ref{tab:parti}.
As Partes I--VI estabelecem o contexto; as Partes VIII--X, as estimativas
necessárias. O passo-chave é um lema técnico de ligação na Parte~VII, de
natureza padrão, cuja verificação fica a cargo do leitor.%
\footnote{O leitor interessado pode consultar o kernel com
\texttt{\#print axioms Disagio.P\_eq\_NP}. A resposta está no log da CI.}
Combinando seis resultados intermediários com o lema de ligação, obtém-se
a tese.
\end{proof}

\begin{table}[ht]
\centering
\begin{tabular}{@{}llr@{}}
\toprule
Parte & Título & Linha \\
\midrule
I    & Preliminares                      & 24 \\
II   & Reduções de sofá                  & 367 \\
III  & A geladeira como oráculo          & 709 \\
IV   & Cook--Levin, mas à noite          & 1057 \\
V    & Encheção de linguiça estratégica  & 1406 \\
VI   & Barreiras superadas               & 1748 \\
VII  & Lemas de ligação                  & 2084 \\
VIII & Ainda preliminares                & 2431 \\
IX   & Estimativas assintóticas          & 2774 \\
X    & Rumo à conclusão                  & 3119 \\
\bottomrule
\end{tabular}
\caption{Estrutura da demonstração. A coluna \emph{Linha} indica onde
começa cada parte em \texttt{PNP.lean}. Ninguém lê além do trecentésimo
lema, razão pela qual escrevemos 780.}
\label{tab:parti}
\end{table}

\begin{osservazione}
A Parte~VIII vem depois da Parte~VII. Trata-se de uma escolha deliberada:
os preliminares, colocados depois do resultado-chave, são mais fáceis de
pular.
\end{osservazione}

\section{Um resultado sem axiomas}

Em benefício dos revisores mais céticos, apresentamos também uma versão do
resultado que não depende de axioma algum, nem mesmo de \texttt{propext}.

\begin{corollario}\label{cor:lett}
Para todos $P, N \in \mathbb{N}$ com $N = 1$, vale $P = N \cdot P$.
\end{corollario}

\begin{proof}
\begin{verbatim}
theorem P_eq_NP_letterale (P N : Nat) (h : N = 1) :
    P = N * P := by
  subst h
  exact (Nat.one_mul P).symm
\end{verbatim}
\end{proof}

O Corolário~\ref{cor:lett} deixa claro que a dificuldade histórica do
problema era, em grande parte, notacional.

\section{Avaliação experimental}

Executamos \texttt{lake build} no Lean~4 v4.14 sem dependências externas.
A compilação termina com sucesso; o selo da CI está verde. O
\texttt{PNP.lean} é gerado de forma determinística pelo \texttt{gen.py} com
semente fixa, o que torna o resultado reprodutível por qualquer pessoa que
tenha um interpretador Python e nenhum pudor.

\section{Conclusões}

Demonstramos que $\mathsf{P} = \mathsf{NP}$. As consequências para a
criptografia são graves e ficam para trabalhos futuros. O prêmio do Clay
Mathematics Institute está aguardando o Pix cair.

\subsection*{Agradecimentos}
Agradecemos à geladeira, aberta e iluminada durante todo o trabalho, e ao
kernel do Lean, que aceitou tudo sem fazer perguntas.

\begin{thebibliography}{9}
\bibitem{aw} S.~Aaronson, A.~Wigderson,
  \emph{Algebrization: A new barrier in complexity theory},
  ACM Trans. Comput. Theory \textbf{1} (2009), n.º~1.
\bibitem{bgs} T.~Baker, J.~Gill, R.~Solovay,
  \emph{Relativizations of the $\mathcal{P} =? \mathcal{NP}$ question},
  SIAM J. Comput. \textbf{4} (1975), 431--442.
\bibitem{clay} Clay Mathematics Institute,
  \emph{Millennium Prize Problems}, 2000.
\bibitem{cook} S.~A.~Cook,
  \emph{The complexity of theorem-proving procedures},
  Proc. 3rd ACM STOC (1971), 151--158.
\bibitem{levin} L.~A.~Levin,
  \emph{Universal sequential search problems},
  Probl. Peredachi Inf. \textbf{9} (1973), n.º~3, 115--116.
\bibitem{rr} A.~A.~Razborov, S.~Rudich,
  \emph{Natural proofs},
  J. Comput. System Sci. \textbf{55} (1997), 24--35.
\bibitem{pnp} Disagio S.A.,
  \emph{PNP.lean}, \url{https://github.com/vjt/pnp}, 2026.
\end{thebibliography}

\end{document}
